\chapter{이 문서에 대하여}

\section{누구를 위한 문서인가}

이 문서는 프로그래밍은 해 보았지만 전기 회로나 FPGA, 디지털 회로 설계는 처음인 사람을 위해 썼다. 목표는
Project Christopher의 기계와 코드를 보고 ``이것이 왜 이렇게 생겼고 어떻게 동작하는가''를 이해하는 것이다.
설치와 실행 방법을 안내하는 문서가 아니다. 그 내용은 부록~\ref{app:run}에 짧게만 있다.

Project Christopher는 제2차 세계대전 때 에니그마 암호를 푸는 데 쓰인 기계 \term{봄베}{Bombe}를
재현하는 프로젝트다. 최종 목표는 드럼이 실제로 돌아가는 기계식 실물이고, 전시해서 관람객이 직접 조작해
볼 수 있는 수준을 설계 요구사항으로 둔다. 전시는 같은 봄베를 세 방식으로 만든 세 대로 구성한다.
기어로 드럼을 돌리는 기어식, 릴레이식, 그리고 움직이는 부품이 없는 전자식(FPGA)이다. 나란히 두면 같은
계산이 기계에서 전기기계로, 다시 전자로 옮겨 온 과정이 보인다. 프로그램(Rust)은 세 대 모두가 내야 할
답을 정하는 기준 모델이다.

\section{이 문서의 줄기}

세 기계는 생김새가 전혀 다르지만 같은 일을 한다. 그 일은 두 가지로 요약된다(\ref{ch:connect}장, \ref{ch:find}장).

\begin{enumerate}
\item 드럼 위치마다 676가닥의 선 사이에 \emph{연결을 만든다}.
\item 전압을 건 선에서 \emph{이어진 선을 모두 찾는다}.
\end{enumerate}

기어식과 릴레이식에서 둘째 일은 전류가 공짜로 해 준다. 두 기계의 차이는 첫째 일, 곧 연결을 드럼의 회전으로
바꾸는가 접점으로 바꾸는가에 있다. FPGA에는 저절로 퍼지는 전류가 없어서 둘째 일을 회로로 직접 지어야 하고,
이것이 FPGA 설계의 대부분을 정한다. 그래서 이 문서는 연결에서 시작해 기어식, 릴레이식, FPGA의 순서로 간다.
뒤의 기계는 앞의 기계가 하던 일을 다른 수단으로 바꾼 것으로 설명한다.

\section{읽는 순서}

부마다 알파벳 하나를 붙였다. 장의 파일 이름도 같은 알파벳으로 시작한다.

\begin{itemize}
\item \textbf{A부: 연결과 탐색.} 봄베가 묻는 것, 그 물음이 왜 ``선과 선의 연결''이 되는지, 이어진 선을 찾는 세 가지
방법. 코드도 회로도 없이 설명한다. 나머지 모든 부가 이 두 장 위에 서므로 먼저 읽는다.
\item \textbf{B부: 기준 답.} Rust로 쓴 소프트웨어 모델이다. 모든 기계가 내야 할 답을 정하고, 봄베의 계산을 가장 읽기
쉬운 형태로 보여 준다.
\item \textbf{C부: 전기와 릴레이.} 실물 두 기계가 함께 쓰는 부품이다. 전압·전류·다이오드와 회로도
읽기(\ref{ch:39-electric-basics}장), 릴레이와 접점으로 논리와 기억을 만드는 법(\ref{ch:relay}장), 느리고 들쭉날쭉한
부품으로 시간을 정하는 법과 테스트 레지스터(\ref{ch:relay-timing}장).
\item \textbf{D부: 기어식.} 드럼이 돌아 연결을 바꾼다. 드럼과 접점, 기어열과 자리올림, 슬롯 원판의 시간(\ref{ch:50-gear-mechanics}장).
\item \textbf{E부: 릴레이식.} 회전을 접점으로 바꾼다. 접점 나무, 이진 시프터, 계수기(\ref{ch:relay-machine}장)와
박자를 치는 순서기(\ref{ch:relay-sequencer}장).
\item \textbf{F부: 릴레이에서 FPGA로.} 릴레이 회로를 FPGA로 옮길 때 무엇이 그대로이고 무엇이 바뀌는지를 개념별로
본다. 회로를 글로 기술한다는 것, 클럭과 레지스터, 자원과 시간, 검증, 그리고 LUT 40만 개짜리 실패 사례.
\item \textbf{G부: FPGA 코드.} 기준 RTL과 실제 칩에 올리는 엔진의 코드를 읽는다.
\item \textbf{H부: 실물로 만들고 검증하기.} 회로도와 기판을 코드로 만드는 법(\ref{ch:boardgen}장), 그 기판을
기준 모델과 대조하는 법(\ref{ch:boardverify}장), 전원과 안전(\ref{ch:44-power-safety}장), 기어식의 파라메트릭
CAD(\ref{ch:51-gear-cad}장)와 제어 보드(\ref{ch:52-gear-control}장).
\item \textbf{I부: 운영 소프트웨어.} 일을 나누어 돌리는 탐색(\ref{ch:14-rust-search}장), 파이썬에서 쓰기(\ref{ch:15-python-api}장),
FPGA 보드와 이야기하는 호스트와 전시 화면(\ref{ch:31-fpga-host}장), 관람객이 쓰는 메뉴 카드(\ref{ch:32-menu-cards}장).
\end{itemize}

봄베의 원리를 이미 알면 B부를 훑기만 해도 된다. 실물 기계에 관심이 있으면 A--E부와 H부를, FPGA에 관심이 있으면
A--B부, C부의 릴레이 장, E부를 읽고 F--G부로 가면 된다. F부는 C부와 E부의 회로를 출발점으로 삼는다.

\section{설명하는 방식}

개념 하나를 세 단계로 설명한다.

\Why 그 개념이 그렇게 될 수밖에 없는 원리와 메커니즘을 먼저 적는다.

\General 이 프로젝트와 무관한 교과서적인 회로나 일상의 예를 든다. 카운터, 깜박이는 LED, 덧셈기, FIFO,
신호 처리 필터 같은 것들이다.

\Ours 같은 개념이 이 프로젝트의 코드와 기계에서 나타나는 곳을 보인다.

코드는 저장소의 실제 코드에서 발췌했다. 지면을 아끼려고 주석 일부와 오류 처리 등을 생략한 곳이 있으며,
생략한 곳은 \code{...}로 표시했다. 코드 안의 주석은 원문대로 영어다.
